\chapter{Énergie~: formes et conservation}\label{ch:g10:energy-conservation}

Une voiture en mouvement, un téléphone chargé, un four chaud et un
marteau levé partagent une propriété~: chacun peut faire arriver quelque
chose. La physique mesure cette capacité avec un seul nombre ---
l'\omterm{def:g10:energy-conservation:energy}{énergie} --- et tient ses comptes sous une règle de fer~: le
total ne change jamais, il ne fait que se déplacer et changer de
costume. Ce chapitre ouvre le grand livre~: les formes que prend
l'\omterm{def:g10:energy-conservation:energy}{énergie}, les chaînes qu'elle parcourt, et ce qu'achète un
\omterm{def:g10:energy-conservation:kwh}{kilowatt-heure}.

\section{Énergie~: la monnaie commune}

\begin{definition}[Énergie et le joule]\label{def:g10:energy-conservation:energy}
L'\emph{énergie}\index{énergie} mesure la capacité d'un système à
produire des changements~: mettre la matière en mouvement, la soulever,
la chauffer, l'éclairer. Son \omterm{def:g10:orders-of-magnitude:unit}{unité} est le
\emph{joule}\index{joule} (\unit{J}). Stockée, transférée, convertie ---
chaque branche de la physique compte dans cette seule monnaie.
\end{definition}

\begin{definition}[Énergie cinétique]\label{def:g10:energy-conservation:kinetic}
Un corps de masse $m$ (en \unit{kg}) se déplaçant à la vitesse $v$ (en
\unit{m/s}) porte l'\emph{énergie cinétique}\index{énergie cinétique}
$E_k = \tfrac12\, m v^2$~: doubler la masse la double, doubler la
vitesse la quadruple.
\end{definition}

\begin{example}[Voiture contre piéton]\label{ex:g10:energy-conservation:kinetic}
Une voiture de \qty{1300}{kg} à \qty{100}{km/h} $= \qty{27.8}{m/s}$
porte $E_k = \tfrac12 \times 1300 \times 27.8^2 \approx \qty{5.0e5}{J}$~;
un piéton de \qty{70}{kg} à \qty{1.5}{m/s},
$\tfrac12 \times 70 \times 1.5^2 \approx \qty{79}{J}$ --- six mille
fois moins. La vitesse est chère~: le carré s'en charge.
\end{example}

\begin{definition}[Énergie potentielle de gravitation]\label{def:g10:energy-conservation:potential}
Un corps de masse $m$ à la hauteur $h$ (en \unit{m}) au-dessus d'un
niveau de référence choisi stocke l'\emph{énergie potentielle de
gravitation}\index{énergie potentielle!de gravitation} $E_p = m g h$,
avec $g = \qty{9.81}{N/kg}$ --- valable près du sol, où $g$ est
effectivement constant. Seules les \emph{différences} de hauteur
comptent~: le niveau où $E_p = 0$ se choisit librement.
\end{definition}

\begin{example}[Château d'eau]\label{ex:g10:energy-conservation:potential}
Chaque \omterm{def:g10:orders-of-magnitude:unit}{mètre} cube (\qty{1000}{kg}) d'eau stocké \qty{30}{m}
en haut d'un château d'eau tient $1000 \times 9.81 \times 30 \approx
\qty{2.9e5}{J}$ --- environ six dixièmes de l'\omterm{def:g10:energy-conservation:kinetic}{énergie
cinétique} de la voiture ci-dessus. Les villes stockent
l'\omterm{def:g10:energy-conservation:energy}{énergie} en hauteur.
\end{example}

\begin{definition}[Les autres formes]\label{def:g10:energy-conservation:forms}
L'\omterm{def:g10:energy-conservation:energy}{énergie} se cache aussi sous des costumes moins visibles~:
\emph{énergie thermique}\index{énergie thermique}, l'agitation
désordonnée des molécules d'un corps (plus chaud, plus il y en a)~;
\emph{énergie chimique}\index{énergie chimique}, stockée dans les
liaisons moléculaires --- nourriture, bois, essence, batteries~;
\emph{énergie électrique}\index{énergie électrique}, portée par les
courants de l'usine à l'utilisateur~; \emph{énergie
radiative}\index{énergie radiative}, portée par la lumière même à
travers le vide --- comment l'\omterm{def:g10:energy-conservation:energy}{énergie} du Soleil nous
atteint~; \emph{énergie nucléaire}\index{énergie nucléaire}, stockée
à l'intérieur des noyaux atomiques, libérée dans les étoiles
(\cref{ch:g12:nuclear-energy}).
\end{definition}

\section{Transferts et chaînes de conversion}

\begin{definition}[Transfert, conversion, chaîne]\label{def:g10:energy-conservation:transfer}
Un \emph{transfert d'énergie}\index{transfert d'énergie} déplace
l'\omterm{def:g10:energy-conservation:energy}{énergie} entre systèmes --- des
\emph{réservoirs}\index{réservoir (d'énergie)}~; une \emph{conversion
d'énergie}\index{conversion d'énergie} change sa forme. Un dispositif
faisant l'un ou l'autre (turbine, moteur, lampe, muscle) est un
\emph{convertisseur}\index{convertisseur}. Une \emph{chaîne de
conversion}\index{chaîne de conversion} figure le voyage~: des boîtes
pour les réservoirs et les formes, des flèches étiquetées pour les
transferts.
\end{definition}

\begin{omfigure}
\begin{tikzpicture}[font=\small,
  box/.style={draw, rounded corners, fill=omDef!10, minimum height=10mm, align=center},
  lab/.style={font=\footnotesize, align=center}]
  \node[box] (res)  at (0,0)   {réservoir d'eau\\(potentielle)};
  \node[box] (tur)  at (4.1,0) {turbine qui tourne\\(cinétique)};
  \node[box] (grid) at (8.2,0) {foyers sur le réseau\\(électrique)};
  \draw[-Stealth, thick, omThm] (res) -- (tur)
    node[midway, above, lab] {chute dans\\la conduite forcée};
  \draw[-Stealth, thick, omThm] (tur) -- (grid)
    node[midway, above, lab] {alternateur};
  \draw[-Stealth, omExo] (tur.south) -- ++(0,-0.7)
    node[below, lab] {chaleur (frottement,\\turbulence)};
  \draw[-Stealth, omExo] (grid.south) -- ++(0,-0.7)
    node[below, lab] {chaleur dans les fils};
\end{tikzpicture}

{\small \omterm{def:g10:energy-conservation:transfer}{Chaîne de conversion} d'un barrage hydroélectrique~: les
boîtes sont des \omterm{def:g10:energy-conservation:transfer}{réservoirs}, les flèches des
transferts --- chaque vraie flèche fuit de la chaleur.}
\end{omfigure}

\begin{method}[Dessiner une chaîne de conversion]\label{met:g10:energy-conservation:chain}
\begin{enumerate}
  \item Identifier le réservoir initial et la forme (qu'est-ce qui est
        consommé~?).
  \item Suivre l'\omterm{def:g10:energy-conservation:energy}{énergie} \omterm{def:g10:energy-conservation:transfer}{convertisseur}
        par \omterm{def:g10:energy-conservation:transfer}{convertisseur}~: une flèche par transfert, la
        forme nommée dans chaque boîte.
  \item Ajouter une flèche de chaleur à chaque vrai
        \omterm{def:g10:energy-conservation:transfer}{convertisseur} --- un peu d'\omterm{def:g10:energy-conservation:energy}{énergie}
        fuit toujours en chaleur (\omterm{def:g10:inertia:inventory}{frottement}, fils,
        échappement).
\end{enumerate}
\end{method}

Un pendule parcourt la chaîne la plus simple de toutes --- toute
potentielle en haut du balancement, toute cinétique au point le plus bas
(en prenant $E_p = 0$ là)~:

\begin{omfigure}
\begin{tikzpicture}[font=\small]
  \draw[-Stealth] (0,0) -- (0,2.5) node[above] {$E$}; \draw (0,0) -- (2.8,0);
  \draw[dashed, omExo] (0,2.0) -- (2.8,2.0) node[right] {total};
  \fill[omDef!60] (0.5,0) rectangle (1.1,2.0); \draw[omThm, very thick] (1.7,0) -- (2.3,0);
  \node at (0.8,-0.4) {$E_p$}; \node at (2.0,-0.4) {$E_k$};
  \node at (1.4,-1.0) {haut du balancement};
\end{tikzpicture}\qquad
\begin{tikzpicture}[font=\small]
  \draw[-Stealth] (0,0) -- (0,2.5) node[above] {$E$}; \draw (0,0) -- (2.8,0);
  \draw[dashed, omExo] (0,2.0) -- (2.8,2.0) node[right] {total};
  \draw[omDef, very thick] (0.5,0) -- (1.1,0); \fill[omThm!60] (1.7,0) rectangle (2.3,2.0);
  \node at (0.8,-0.4) {$E_p$}; \node at (2.0,-0.4) {$E_k$};
  \node at (1.4,-1.0) {point le plus bas};
\end{tikzpicture}

{\small Le budget du pendule en deux positions~: $E_p$ et $E_k$
s'échangent tandis que leur somme (en tirets) reste en place.}
\end{omfigure}

\section{Conservation de l'énergie}

\begin{theorem}[Conservation de l'énergie]\label{thm:g10:energy-conservation:conservation}
\index{conservation de l'énergie}
Le total d'\omterm{def:g10:energy-conservation:energy}{énergie} d'un système isolé --- un qui n'échange
rien avec l'extérieur --- est constant. L'\omterm{def:g10:energy-conservation:energy}{énergie} n'est
jamais créée et jamais détruite~: elle ne change que de forme et de
lieu, et ce qu'un réservoir perd, les autres le gagnent,
\omterm{def:g10:energy-conservation:energy}{joule} pour \omterm{def:g10:energy-conservation:energy}{joule}.
\end{theorem}
\admitted

\begin{remark}\label{rem:g10:energy-conservation:admitted}
Aucune expérience n'a jamais pris ce principe en défaut~; ici on le
prend comme donné et on l'utilise numériquement.
Le~\cref{ch:g12:work-and-energy} dérive sa partie mécanique
($E_k + E_p$) des lois de Newton~; la comptabilité complète attend les
volumes de licence.
\end{remark}

\begin{example}[Chute libre, tarifée en joules]\label{ex:g10:energy-conservation:fall}
Une balle de \qty{2.0}{kg} tombe de \qty{10}{m} (résistance de l'air
négligeable)~: elle part avec
$E_p = 2.0 \times 9.81 \times 10 \approx \qty{196}{J}$, arrive avec
$E_k = \qty{196}{J}$, donc $v = \sqrt{2 \times 196/2.0} =
\qty{14}{m/s} \approx \qty{50}{km/h}$. Aucune \omterm{def:g10:inertia:force}{force},
aucun chronomètre~: la conservation seule tarife l'impact.
\end{example}

\section{Rendement}

\begin{definition}[Rendement]\label{def:g10:energy-conservation:efficiency}
Un \omterm{def:g10:energy-conservation:transfer}{convertisseur} consomme $E_{\text{consumed}}$ et livre
la forme voulue $E_{\text{useful}}$~; son
\emph{rendement}\index{rendement} est
\[
\eta = \frac{E_{\text{useful}}}{E_{\text{consumed}}} < 1 .
\]
La conservation interdit $\eta > 1$, et les vrais
\omterm{def:g10:energy-conservation:transfer}{convertisseurs} n'atteignent jamais $1$~: l'\omterm{def:g10:energy-conservation:energy}{énergie}
manquante n'est pas détruite --- elle part en chaleur.
\end{definition}

\begin{example}[Où va l'essence]\label{ex:g10:energy-conservation:engine}
Un \omterm{def:g10:orders-of-magnitude:unit}{kilogramme} d'essence stocke \qty{45}{MJ}
d'\omterm{def:g10:energy-conservation:forms}{énergie chimique}~; un moteur de voiture le convertit avec
$\eta \approx 0.30$~: \qty{13.5}{MJ} de mouvement, \qty{31.5}{MJ}
chauffant l'échappement, le radiateur et la rue. Les deux tiers de
chaque réservoir réchauffent l'atmosphère.
\end{example}

\begin{omfigure}
\begin{tikzpicture}[font=\small]
  \fill[omDef!40] (0,0) rectangle (3,1.5);
  \node[align=center] at (1.5,0.75) {essence\\\qty{45}{MJ}};
  \fill[omMeth!60] (3,1.05) -- (5.6,1.55) -- (5.6,2.0) -- (3,1.5);
  \fill[omMeth!60] (5.6,1.45) -- (6.3,1.775) -- (5.6,2.1);
  \node[right] at (6.35,1.78) {mouvement~: \qty{13.5}{MJ} ($30\%$)};
  \fill[omExo!70] (3,0) -- (5.6,-1.0) -- (5.6,0.05) -- (3,1.05);
  \fill[omExo!70] (5.6,-1.1) -- (6.3,-0.475) -- (5.6,0.15);
  \node[right] at (6.35,-0.48) {chaleur~: \qty{31.5}{MJ} ($70\%$)};
\end{tikzpicture}

{\small Le budget du moteur de voiture, largeurs de flèches à l'échelle~:
sur \qty{45}{MJ} consommés, \qty{13.5}{MJ} font avancer la voiture~; le
reste est chaleur.}
\end{omfigure}

\begin{remark}[Les rendements se multiplient]\label{rem:g10:energy-conservation:chain-eta}
Des \omterm{def:g10:energy-conservation:transfer}{convertisseurs} en chaîne multiplient leurs
rendements~: une centrale (\omterm{def:g10:energy-conservation:power}{puissance} $0{,}38$) alimentant
le réseau ($0{,}94$) alimentant un chargeur ($0{,}85$) livre
$\eta = 0.38 \times 0.94 \times 0.85 \approx 0.30$. Les longues chaînes
sont des tuyaux qui fuient~: chaque flèche taxe le flux.
\end{remark}

\section{Puissance, le watt et le kilowatt-heure}

\begin{definition}[Puissance]\label{def:g10:energy-conservation:power}
La \emph{puissance}\index{puissance} est le débit auquel
l'\omterm{def:g10:energy-conservation:energy}{énergie} s'écoule, $P = E/t$, mesurée en
\emph{watts}\index{watt} (\unit{W}), avec
$\qty{1}{W} = \qty{1}{J/s}$. L'\omterm{def:g10:energy-conservation:energy}{énergie} est une
quantité, la puissance un flux --- une baignoire face à son robinet.
\end{definition}

\begin{example}[Bouilloire contre humain]\label{ex:g10:energy-conservation:power}
Une bouilloire de \qty{2200}{W} tournant \qty{150}{s} utilise
$E = 2200 \times 150 = \qty{3.3e5}{J}$. Un humain tourne sur environ
\qty{10}{MJ} de nourriture par jour, une moyenne
$\num{1.0e7}/\num{86400} \approx \qty{116}{W}$~: vous tournez au
ralenti comme une vieille ampoule brillante.
\end{example}

\begin{definition}[Le kilowatt-heure]\label{def:g10:energy-conservation:kwh}
Le \emph{kilowatt-heure}\index{kilowatt-heure} (\unit{kW.h}) est
l'\omterm{def:g10:energy-conservation:energy}{énergie} d'un flux de \qty{1}{kW} tournant une heure~:
$\qty{1}{kW.h} = 1000 \times 3600\,\unit{J} = \qty{3.6}{MJ}$ ---
l'\omterm{def:g10:orders-of-magnitude:unit}{unité} que comptent les compteurs d'électricité~:
\omterm{def:g10:energy-conservation:power}{puissance} en \unit{kW} multipliée par des heures.
\end{definition}

\begin{example}[Lire la facture]\label{ex:g10:energy-conservation:bill}
Une facture est une soustraction et une multiplication~: nouvelle
lecture du compteur moins l'ancienne donne les
\omterm{def:g10:energy-conservation:kwh}{kilowatt-heures}~; multipliés par le prix de
l'\omterm{def:g10:orders-of-magnitude:unit}{unité} (environ $0{,}25$~euro par \unit{kW.h}), plus un
abonnement fixe, donne le total. Un radiateur de \qty{2}{kW} tournant
\qty{2.5}{h} chaque nuit utilise \qty{5}{kW.h} par jour~: $38$~euros
par mois à lui seul.
\end{example}

\begin{example}[Ordres de grandeur]\label{ex:g10:energy-conservation:orders}
À connaître par cœur~: soulever une pomme d'un
\omterm{def:g10:orders-of-magnitude:unit}{mètre}, environ \qty{1}{J}~; une batterie de téléphone,
$\qty{10}{W.h} \approx \qty{4e4}{J}$~; un bol de céréales,
\qty{1.5}{MJ}~; votre nourriture pour un jour, \qty{10}{MJ}~; un
\omterm{def:g10:orders-of-magnitude:unit}{kilogramme} d'essence, \qty{45}{MJ} --- presque exactement
l'électricité quotidienne d'un foyer ($\approx \qty{12}{kW.h}$)~; un
éclair, \qty{5}{GJ}~: cent jours-foyer, livrés en un flash.
\end{example}

\begin{omfigure}
\begin{tikzpicture}[font=\small, y=0.5cm]
  \draw[-Stealth] (0,-0.5) -- (0,10.8) node[above] {$E$ (\unit{J})};
  \foreach \n in {0,2,4,6,8,10}
    {\draw (-0.12,\n) -- (0.12,\n); \node[left, font=\footnotesize] at (-0.15,\n) {$10^{\n}$};}
  \fill[omDef] (0,0) circle (1.6pt) node[right=4pt] {pomme soulevée de \qty{1}{m}};
  \fill[omDef] (0,4.56) circle (1.6pt) node[right=4pt] {batterie de téléphone (\qty{10}{W.h})};
  \fill[omDef] (0,6.18) circle (1.6pt) node[right=4pt] {bol de céréales};
  \fill[omDef] (0,7) circle (1.6pt) node[right=4pt] {votre nourriture, un jour (\qty{10}{MJ})};
  \fill[omThm] (0,7.65) circle (1.6pt) node[right=4pt] {\qty{1}{kg} d'essence $\approx$ jour de foyer};
  \fill[omThm] (0,9.7) circle (1.6pt) node[right=4pt] {éclair ($\approx\qty{5}{GJ}$)};
\end{tikzpicture}

{\small Une échelle d'\omterm{def:g10:energy-conservation:energy}{énergie}~: chaque barreau est un facteur
$100$. La vie quotidienne couvre dix ordres de grandeur.}
\end{omfigure}

\section{Exercices}

\begin{exercise}[$\star$]\label{exo:g10:energy-conservation:1}
Une voiture de \qty{1200}{kg} roule à \qty{90}{km/h}, c'est-à-dire
\qty{25}{m/s}~: calculer son \omterm{def:g10:energy-conservation:kinetic}{énergie cinétique}. La
refaire à \qty{130}{km/h} (\qty{36.1}{m/s})~: que $44\%$ de vitesse en
plus a-t-elle fait à $E_k$~?
\end{exercise}

\begin{exercise}[$\star$]\label{exo:g10:energy-conservation:2}
Un randonneur de masse \qty{72}{kg} grimpe de l'altitude \qty{1200}{m}
à \qty{1650}{m}. Calculer l'\omterm{def:g10:energy-conservation:energy}{énergie} potentielle gagnée,
la convertir en \omterm{def:g10:energy-conservation:kwh}{kilowatt-heures}, et la tarifer à
$0{,}25$~euro par \unit{kW.h}.
\end{exercise}

\begin{exercise}[$\star$]\label{exo:g10:energy-conservation:3}
Convertir~: \qty{1}{kW.h} en \omterm{def:g10:energy-conservation:energy}{joules}~; une batterie de
téléphone de \qty{12}{W.h} en kilojoules~; \qty{45}{MJ} (un
\omterm{def:g10:orders-of-magnitude:unit}{kilogramme} d'essence) en
\omterm{def:g10:energy-conservation:kwh}{kilowatt-heures}~; votre journée de nourriture de
\qty{10}{MJ} en \omterm{def:g10:energy-conservation:kwh}{kilowatt-heures}.
\end{exercise}

\begin{exercise}[$\star$]\label{exo:g10:energy-conservation:4}
Une bouilloire de \qty{2200}{W} tourne \qty{150}{s}. Calculer
l'\omterm{def:g10:energy-conservation:energy}{énergie} utilisée, en \omterm{def:g10:energy-conservation:energy}{joules} et
en \omterm{def:g10:energy-conservation:kwh}{kilowatt-heures}. Puis calculer la
\omterm{def:g10:energy-conservation:power}{puissance} moyenne d'un humain tournant sur
\qty{10}{MJ} par jour.
\end{exercise}

\begin{exercise}[$\star$]\label{exo:g10:energy-conservation:5}
Écrire la \omterm{def:g10:energy-conservation:transfer}{chaîne de conversion}
(\cref{met:g10:energy-conservation:chain}) de~: une dynamo de vélo
alimentant une lampe~; un chauffe-eau à gaz~; un panneau solaire
chargeant un téléphone. Marquer les fuites de chaleur.
\end{exercise}

\begin{exercise}[$\star\star$]\label{exo:g10:energy-conservation:6}
Un moteur brûle \qty{2.0}{kg} d'essence (\qty{45}{MJ/kg}) et livre
\qty{27}{MJ} de mouvement. Calculer son
\omterm{def:g10:energy-conservation:efficiency}{rendement} et l'\omterm{def:g10:energy-conservation:energy}{énergie}
libérée en chaleur. Où va cette chaleur~?
\end{exercise}

\begin{exercise}[$\star\star$]\label{exo:g10:energy-conservation:7}
Un plongeur de falaise de \qty{70}{kg} quitte une falaise de
\qty{20}{m} (résistance de l'air négligeable). En utilisant la
conservation, calculer l'\omterm{def:g10:energy-conservation:kinetic}{énergie cinétique} à l'eau
et la vitesse d'entrée, en \unit{m/s} et \unit{km/h}.
\end{exercise}

\begin{exercise}[$\star\star$]\label{exo:g10:energy-conservation:8}
Un compteur a lu \qty{41236}{kW.h} le 1\textsuperscript{er} mars,
\qty{41562}{kW.h} $30$~jours plus tard~:
\omterm{def:g10:energy-conservation:energy}{énergie} utilisée, moyenne journalière, facture à
$0{,}26$~euro par \unit{kW.h} plus $12$~euros fixes, et
\omterm{def:g10:energy-conservation:power}{puissance} moyenne en \omterm{def:g10:energy-conservation:power}{watts}~?
\end{exercise}

\begin{exercise}[$\star\star$]\label{exo:g10:energy-conservation:9}
Une vieille ampoule de \qty{60}{W} ne transforme que $5\%$ de sa
consommation en lumière~; une LED donne la même lumière avec
\qty{9}{W}. Calculer la \omterm{def:g10:energy-conservation:power}{puissance} lumineuse de
l'ampoule et le \omterm{def:g10:energy-conservation:efficiency}{rendement} de la LED~; puis, à
\qty{4}{h} par jour, la consommation annuelle de chacune et l'économie
à $0{,}25$~euro par \unit{kW.h}.
\end{exercise}

\begin{exercise}[$\star\star$]\label{exo:g10:energy-conservation:10}
Un barrage laisse passer \qty{150}{m^3} d'eau (\qty{1.5e5}{kg}) par
seconde sur une chute de \qty{40}{m}. Calculer l'\omterm{def:g10:energy-conservation:energy}{énergie}
potentielle libérée chaque seconde --- la
\omterm{def:g10:energy-conservation:power}{puissance} disponible --- puis la
\omterm{def:g10:energy-conservation:power}{puissance} électrique à $\eta = 0.90$, et combien de
foyers de moyenne \qty{525}{W} elle alimente.
\end{exercise}

\begin{exercise}[$\star\star$]\label{exo:g10:energy-conservation:11}
Un pendule de \qty{0.20}{kg} est lâché \qty{5.0}{cm} au-dessus de son
point le plus bas~: calculer sa vitesse là. Après quelques minutes il
pend immobile~: où est passée son \omterm{def:g10:energy-conservation:energy}{énergie}, et pourquoi
cela ne contredit-il pas
le~\cref{thm:g10:energy-conservation:conservation}~?
\end{exercise}

\begin{exercise}[$\star\star\star$]\label{exo:g10:energy-conservation:12}
Une centrale à charbon ($\eta_1 = 0.38$), le réseau
($\eta_2 = 0.94$) et un chargeur ($\eta_3 = 0.85$) rechargent une
batterie de téléphone de \qty{12}{W.h}. Calculer le
\omterm{def:g10:energy-conservation:efficiency}{rendement} global~; l'\omterm{def:g10:energy-conservation:forms}{énergie
chimique} consommée à la centrale, en \unit{W.h} et kilojoules~; et le
charbon (\qty{25}{MJ/kg}) par charge, puis par an de charges
quotidiennes.
\end{exercise}

\begin{exercise}[$\star\star\star$]\label{exo:g10:energy-conservation:13}
Un éclair porte environ \qty{5}{GJ}. Convertir en
\omterm{def:g10:energy-conservation:kwh}{kilowatt-heures} et en jours d'un foyer de
\qty{12.6}{kW.h} par jour. Il est livré en environ
\qty{100}{\micro s}~: calculer la \omterm{def:g10:energy-conservation:power}{puissance}, la
comparer à la \omterm{def:g10:energy-conservation:power}{puissance} électrique moyenne de
l'humanité (\qty{3}{TW}), et expliquer pourquoi personne ne récolte la
foudre.
\end{exercise}

\begin{exercise}[$\star\star\star$]\label{exo:g10:energy-conservation:14}
Une voiture de \qty{1300}{kg} freine de \qty{110}{km/h}
(\qty{30.6}{m/s}) jusqu'à l'arrêt. Calculer l'\omterm{def:g10:energy-conservation:kinetic}{énergie
cinétique} dissipée~; quelle forme prend-elle, et où~? Une voiture
électrique en récupère $60\%$ dans sa batterie~:
\omterm{def:g10:energy-conservation:energy}{énergie} par arrêt, et arrêts nécessaires pour
remplir un \omterm{def:g10:energy-conservation:kwh}{kilowatt-heure}~?
\end{exercise}

\begin{exercise}[$\star\star\star$]\label{exo:g10:energy-conservation:15}
Un cycliste et son vélo (\qty{85}{kg} en tout) grimpent un col,
gagnant \qty{1200}{m} d'altitude~; les muscles convertissent la
nourriture avec $\eta \approx 0.25$. Calculer l'\omterm{def:g10:energy-conservation:energy}{énergie}
mécanique de la montée, l'\omterm{def:g10:energy-conservation:energy}{énergie} alimentaire qu'elle
exige, l'équivalent en bols de céréales de \qty{1.5}{MJ}, et la masse
d'essence (\qty{45}{MJ/kg}) stockant la même
\omterm{def:g10:energy-conservation:energy}{énergie}. Commenter.
\end{exercise}

\section{Problème~: Alimenter une maison pour un jour}

\begin{problem}[{Problème du week-end --- du bol de céréales aux
panneaux du toit~: un audit énergétique complet d'un jour ordinaire
chez soi, et ce qu'un joule coûte selon qui le vend}]
\label{pb:g10:energy-conservation:1}
Un mardi ordinaire, une famille se met en quête de savoir où va son
électricité, armée de ce chapitre, du compteur du couloir et du dos de
la facture. Le verdict de ce soir~: leur jour égale un
\omterm{def:g10:orders-of-magnitude:unit}{kilogramme} d'essence, seize
\omterm{def:g10:orders-of-magnitude:unit}{mètres} carrés de soleil --- ou cinq cyclistes pédalant
jour et nuit.

\medskip
\noindent\textbf{Partie I --- L'audit.}
L'inventaire du jour~: frigo \qty{100}{W} pendant \qty{24}{h}~;
chauffe-eau \qty{2400}{W} pendant \qty{2.5}{h}~; four \qty{2000}{W}
pendant \qty{1.0}{h}~; machine à laver \qty{500}{W} pendant
\qty{2.0}{h}~; éclairage \qty{60}{W} pendant \qty{5.0}{h}~; écrans
\qty{150}{W} pendant \qty{6.0}{h}.
\begin{enumerate}
  \item Calculer l'\omterm{def:g10:energy-conservation:energy}{énergie} de chaque appareil pour le
        jour, en \unit{kW.h}.
  \item Totaliser le jour en \omterm{def:g10:energy-conservation:kwh}{kilowatt-heures}, puis en
        mégajoules~; comparer à un \omterm{def:g10:orders-of-magnitude:unit}{kilogramme} d'essence
        (\qty{45}{MJ}).
  \item Calculer la \omterm{def:g10:energy-conservation:power}{puissance} moyenne de la maison
        sur les \qty{24}{h}~; quelle fraction d'une bouilloire en
        marche est-ce~?
  \item Quel appareil unique domine~? Recalculer le total journalier
        si le chauffe-eau n'a tourné que \qty{1.25}{h}.
  \item Les voyants de veille tirent un permanent \qty{15}{W}~:
        \omterm{def:g10:energy-conservation:energy}{énergie} journalière, part du total, coût
        annuel à $0{,}25$~euro par \unit{kW.h}~?
\end{enumerate}

\medskip
\noindent\textbf{Partie II --- La facture.}
\begin{enumerate}[resume]
  \item En $30$~jours le compteur est passé de \qty{56214}{kW.h} à
        \qty{56604}{kW.h}~: consommation, et moyenne journalière
        confrontée à l'audit~?
  \item Calculer la facture~: $0{,}25$~euro par \unit{kW.h} plus un
        abonnement fixe de $14$~euros.
  \item Prix d'un mégajoule du réseau~? De l'essence, à $1{,}80$~euro
        le litre stockant \qty{34}{MJ}~?
  \item Un bol de céréales de \qty{1.5}{MJ} coûte environ
        $0{,}50$~euro~: son prix par mégajoule~? Classer les trois
        vendeurs de \omterm{def:g10:energy-conservation:energy}{joules}.
\end{enumerate}

\medskip
\noindent\textbf{Partie III --- Le toit.}
\begin{enumerate}[resume]
  \item Le plein soleil livre environ \qty{1000}{W} par
        \omterm{def:g10:orders-of-magnitude:unit}{mètre} carré~; les panneaux le convertissent
        avec $\eta = 0.20$~: \omterm{def:g10:energy-conservation:power}{puissance} électrique de
        \qty{1}{m^2} en plein soleil~?
  \item La région moyenne l'équivalent de \qty{4.0}{h} de plein soleil
        par jour~: rendement journalier de \qty{1}{m^2}~?
  \item Quelle surface de panneaux couvre le jour de
        \qty{12.6}{kW.h} de la famille~? Le toit sud offre
        \qty{20}{m^2}~: cela tient-il~?
  \item En décembre l'équivalent tombe à \qty{1.5}{h}~: la production
        des panneaux, sa part du besoin, et qu'est-ce qui comble
        l'écart~?
  \item La fabrication des panneaux a coûté environ \qty{700}{kW.h} par
        \omterm{def:g10:orders-of-magnitude:unit}{mètre} carré~: temps de retour énergétique, en
        utilisant la question~11~?
\end{enumerate}

\medskip
\noindent\textbf{Partie IV --- Le bol de céréales.}
\begin{enumerate}[resume]
  \item Un vélo-générateur livre \qty{100}{W} d'électricité~: combien
        d'heures de pédalage font \qty{12.6}{kW.h} --- combien de
        personnes pédalant \qty{24}{h} sans arrêt~?
  \item Les muscles tournent à $\eta \approx 0.25$~: l'\omterm{def:g10:energy-conservation:energy}{énergie}
        alimentaire derrière ces \qty{12.6}{kW.h}, en mégajoules et en
        bols de céréales~?
  \item Coût du jour pédalé (bols à $0{,}50$~euro), face à ce que le
        réseau facture pour le même jour~?
  \item Un éclair de \qty{5}{GJ} égale combien de jours-maison~? Compte
        tenu de sa durée de \qty{100}{\micro s}, pourquoi ne peut-il
        quand même pas \omterm{def:g10:energy-conservation:power}{alimenter} la maison~?
  \item Les joules ne se perdent jamais, pourtant les
        \omterm{def:g10:energy-conservation:energy}{joules} du réseau, de l'essence et de la
        nourriture se vendent à des prix très différents~: que
        paie-t-on réellement~?
  \item Final~: énoncer le verdict du jour en une phrase --- un
        \omterm{def:g10:orders-of-magnitude:unit}{kilogramme} d'essence, seize
        \omterm{def:g10:orders-of-magnitude:unit}{mètres} carrés de toit ensoleillé ou cinq cyclistes
        jour et nuit --- et nommer les deux grandeurs (un stock, un
        flux) à ne plus jamais confondre.
\end{enumerate}
\end{problem}
